% GENERATED FILE. Edit canonical lessons, then run npm run generate:books.
% canonical-source-hash: fnv1a64:193aeb3fa3796a6a
% canonical-lessons: ES-C51-repaso-flecha

\chapter{Review --- One Arrow}
\label{ch:review-por-para}
\hlchaptermodality{157}

\begin{chapteropening}
\textbf{By the end of this chapter:} \emph{I can state both directions and place my old por phrases under them.}

Complete ES-C51-repaso-flecha: give both directions, name three old phrases containing por, and say what it means when a sentence takes either.
\end{chapteropening}

\section[Practice]{Review — one arrow}

\label{lesson:ES-C51-repaso-flecha}



\begin{quote}
Nothing new. One question before the table: which way does each preposition point?
\end{quote}

\begin{grammarlens}[title={What you've built}]
\noindent
\begin{tabularx}{\linewidth}{@{}>{\raggedright\arraybackslash}X>{\raggedright\arraybackslash}X>{\raggedright\arraybackslash}X@{}}
\toprule
\textbf{} & \textbf{\textbf{para} — toward} & \textbf{\textbf{por} — through, because} \\
\midrule
built from & \emph{ad}, \textquotedblleft{}toward\textquotedblright{} & \emph{per}, \textquotedblleft{}through\textquotedblright{} \\
the arrow & points \textbf{at} it & comes \textbf{from} it \\
a person & \emph{para ti} — you get it & \emph{por ti} — you caused it \\
a time & \emph{para mañana} — by then & — \\
thanks & — & \emph{gracias por} — for what caused it \\
asking & — & \emph{por favor} — through a favour \\
asking why & — & \emph{¿por qué?} — through what? \\
\bottomrule
\end{tabularx}

Look at the right-hand column. Four of those six \emph{por} entries are phrases you have owned since the opening chapters. This review is not mostly teaching you new Spanish — it is \textbf{re-filing old Spanish under a heading that explains it}.

That happens more than a course usually admits. A great deal of what looks like new grammar late in a book is a name for something you have been doing correctly and blindly for a long time.
\end{grammarlens}

\begin{grammarlens}[title={Grammar Lens: what to do when you are unsure}]
You will still hesitate sometimes. Two things help, in this order:

1. \textbf{Ask what the sentence is about.} A destination, a deadline, a recipient, a purpose — \emph{para}. A cause, a route, a swap, a duration — \emph{por}. 2. \textbf{If it is genuinely both, it is genuinely both.} \emph{Lo hago para ti} and \emph{lo hago por ti} are both good Spanish; you are not choosing right from wrong, you are choosing which thing to say.

That second point matters more than it looks. Most \emph{por}/\emph{para} anxiety comes from believing every sentence has one correct answer waiting to be found. Plenty of sentences take either, and the choice is the meaning.
\end{grammarlens}

\subsection*{Your turn}
\begin{itemize}
\raggedright
  \item \emph{Say it:} \textquotedblleft{}para ti / por ti\textquotedblright{}, \textquotedblleft{}para mañana\textquotedblright{}, \textquotedblleft{}gracias por la comida\textquotedblright{}
\end{itemize}

\emph{Twice through:}

\begin{tcolorbox}[breakable,colback=teal!4,colframe=teal!35!black,title={Before you move on}]
\emph{Para} — which way? (\textbf{At} it.) \emph{Por} — which way? (\textbf{From} it.) And when a sentence takes either? (\textbf{Both are right} — the choice is the meaning.) Next: the two of them in a conversation, where the choice does real work.
\end{tcolorbox}
